%! Problems with Sampling — Unit 3, Lesson 3.3 — Homework (Answer Key)
% MIRROR of ../homework/main.tex: -key for -boxes, the header, and the answers. Nothing else differs.
% NO teachernote here — teacher prose lives in the lesson plan.
% GRADUAL RELEASE: this is the lesson's SCORED practice and the LAST component in the packet.
% Paper homework is the DEFAULT for this course; DeltaMath is an occasional override that
% replaces this page. Started in the last eight minutes of the period and finished after two
% study halls. The AP Practice component before it is extra credit, scored 1 to 5.
% THIRD CONTEXT: the notes teach on one-line surveys and a game forum, Guided Practice runs on a
% university work-hours survey, AP Practice on transit, speed limits, and sleep; these items are
% on a state DOT survey of teen drivers who use a phone while driving.
% Part A: A1 voluntary response with direction, A2 response bias in a RANDOM sample (a trooper
% asks), A3 nonresponse with how the non-answerers differ. Part B: B1 the AP-style multiple
% choice (which change reduces A2's bias) with a required justification, B2 THE CRUX restated —
% "3,000 responses, three times as many, so Plan 1 is more trustworthy". 10 points.
% \answerspace{H}{} reserves exactly H; each \gradedblank inside it is BYTE-IDENTICAL in the key,
% where key mode prints its display. Keep every \boxguard identical in the blank and the key.
\documentclass[12pt]{article}
\usepackage{apstats-article}
\usepackage{apstats-key}   % pulls in -boxes; do NOT also load -boxes
\usepackage{grading}
\gradingassignment{apstats-u03-l03-hw}
\gradingcourse{advanced_statistics}
\newcommand{\answerspace}[2]{\par\nopagebreak\noindent\begin{minipage}[t][#1][t]{\linewidth}%
  \color{keyred}\bfseries #2\end{minipage}\par}

\begin{document}

\pageheader{Unit 3, Lesson 3.3}{Homework --- Answer Key}
{\small\itshape Write each final answer in its box. Only what is inside a box is graded.}\par
% NAMESTRIP: no \namedateperiod here — the name row lives on the cover only.

\vspace{0.04in}

% Authored BYTE-IDENTICALLY in homework/ and homework_key/.
\boxguard[8]
\begin{remindbox}
\small
\textbf{This is your graded homework.} It is scored 1 to 5 and due the first class after two study
halls. Start it now; ask about anything that stalls you before you leave. For every bias you
name, say \textbf{who} is over- or under-represented and whether the estimate is
\textbf{too high or too low}.
\end{remindbox}

\vspace{0.04in}

\boxguard[14]
\begin{notesbox}{Part A --- Teen Drivers and Phones}
\interlinepenalty=10000 % keep each prompt whole when the box breaks
\small
The state Department of Transportation (DOT) wants to estimate the percent of the state's
licensed teen drivers who use a phone while driving. \textbf{Plan 1:} the DOT posts a poll on
its social media page; 3{,}000 teens answer, and \textbf{22\%} say they use a phone while
driving. \textbf{Plan 2:} the DOT selects 1{,}000 teens at random from the state's list of
licensed teen drivers, and a state trooper phones each one to ask the same question.
\begin{enumerate}[label=\textbf{A\arabic*.}, leftmargin=2.2em, itemsep=2pt]
  \item Plan 1: name the bias that comes from \emph{who chose} to answer. Is 22\% likely too
        high or too low? Explain in context.
        \answerspace{2.1in}{\normalcolor\mdseries\hspace*{10.5pt}\gradedblank[id=hw1, type=rubric, points=2, display={Voluntary response bias. Teens who follow the DOT and choose to answer are likely the most safety-minded, so phone users are under-represented: 22\% is likely too low.}]{\linewidth-10.5pt}{0.95cm}}
  \item Plan 2's sample is random, yet its result may still be biased. Name the bias, give its
        direction, and explain in context.
        \answerspace{2.1in}{\normalcolor\mdseries\hspace*{10.5pt}\gradedblank[id=hw2, type=rubric, points=2, display={Response bias. Using a phone while driving is against the law, and teens are unlikely to admit it to a state trooper, so the estimate is likely too low.}]{\linewidth-10.5pt}{0.95cm}}
  \item Only 610 of the 1{,}000 teens in Plan 2 pick up the call. Name this bias, and explain
        how the 390 who did not answer might differ from the 610 who did.
        \answerspace{2.1in}{\normalcolor\mdseries\hspace*{10.5pt}\gradedblank[id=hw3, type=rubric, points=2, display={Nonresponse bias. Teens who let an unknown call ring --- often because they are busy or driving --- may be more likely to use a phone at the wheel, so the estimate is likely too low.}]{\linewidth-10.5pt}{0.95cm}}
\end{enumerate}
\end{notesbox}

\vspace{0.04in}

\boxguard[30]
\begin{notesbox}{Part B --- Fixing the Survey}
\interlinepenalty=10000 % keep each prompt whole when the box breaks
\small
\begin{enumerate}[label=\textbf{B\arabic*.}, leftmargin=2.2em, itemsep=2pt]
  \item \textbf{AP-style.} Which change to Plan 2 would do the most to reduce the bias you
        named in A2? Write its letter in the Letter box, then justify it in a sentence.
        \begin{enumerate}[label=\textbf{(\Alph*)}, leftmargin=2.2em, itemsep=0pt]
          \item Call 5{,}000 teens instead of 1{,}000.
          \item \ans{Let each selected teen answer on an anonymous online form with no name attached.}
          \item Ask, ``Do you put lives at risk by using your phone while driving?''
          \item Call only teens who have completed a driver's education course.
          \item Ask the question in person during each teen's road test.
                \hfill\gradedright{\textit{Letter:}~\gradedblank[id=hw4, type=choice, points=1, answer=B]{0.45in}{0.24in}}\rule[-0.2in]{0pt}{0.48in}
        \end{enumerate}
        \answerspace{2.1in}{\normalcolor\mdseries\hspace*{10.5pt}\gradedblank[id=hw4j, type=rubric, points=1, display={(B): with no name and no trooper, a teen has no reason to hide phone use, so answers are more honest. A bigger sample (A) does not fix response bias.}]{\linewidth-10.5pt}{0.7cm}}
  \item A reporter writes, ``Plan 1 had 3{,}000 responses --- three times as many as Plan 2
        --- so Plan 1 gives the more trustworthy estimate.'' Respond.
        \answerspace{2.1in}{\normalcolor\mdseries\hspace*{10.5pt}\gradedblank[id=hw5, type=rubric, points=2, display={Disagree. Plan 1 is a voluntary response sample, and 3,000 volunteers are still biased. A larger sample reduces variability, not bias, so size cannot fix how Plan 1 got its teens.}]{\linewidth-10.5pt}{0.95cm}}
\end{enumerate}
\end{notesbox}

\vspace{0.04in}

% The next-lesson preview closes the packet, so it lives here rather than in AP Practice.
% A one-line, unbreakable box: guard it by its own height so it moves whole.
\boxguard[3]
\begin{spiralbox}
\small
\textbf{Next: Designing Experiments} --- from watching to doing: treatments, random
assignment, control, and why only an experiment can show cause.
\end{spiralbox}

\end{document}
